In this work, we investigated the impact of epistemic uncertainty on defect detection in medical device manufacturing, a first step particularly important when incorrect predictions can propagate through subsequent root-cause analysis and risk assessment. We showed that explicitly accounting for this uncertainty through abstention provides a simple yet principled mechanism for reducing the risk associated with unreliable model decisions.

Building on a cost-sensitive classification framework, we analyzed a rejection strategy that allows the model to abstain from uncertain predictions rather than forcing a decision for every input. Our experiments on synthetic benchmarks show that abstaining on a relatively small fraction of predictions can reduce decision risk while improving the reliability of the remaining decisions. We further evaluated the approach using real-world medical device report data, demonstrating its applicability beyond controlled synthetic settings. To support reproducibility and further research, we have shared a link to our open-source code.

More broadly, our results suggest that predictive performance alone may not adequately characterize model reliability in high-stakes defect detection. Providing models with the ability to identify and defer uncertain decisions offers a practical mechanism for incorporating uncertainty into automated quality assurance and provides a foundation for integrating such decisions with downstream root-cause analysis, risk assessment, and human review.

